\documentclass[11pt,a4paper]{article}

\usepackage[utf8]{inputenc}
\usepackage[T1]{fontenc}
\usepackage[provide=*,italian]{babel}
\usepackage{lmodern}
\usepackage{microtype}
\usepackage{geometry}
\usepackage{amsmath,amssymb}
\usepackage{booktabs}
\usepackage{longtable}
\usepackage{tabularx}
\usepackage{array}
\usepackage{xcolor}
\usepackage{enumitem}
\usepackage{graphicx}
\usepackage{hyperref}
\usepackage{fancyhdr}
\usepackage{listings}

\geometry{margin=2.5cm}
\hypersetup{
    colorlinks=true,
    linkcolor=blue!55!black,
    citecolor=blue!55!black,
    urlcolor=blue!55!black,
    pdftitle={Registrazione learning-based del phantom cervicale},
    pdfauthor={Giovanni Contursi}
}

\pagestyle{fancy}
\fancyhf{}
\lhead{Medical Robotics}
\rhead{Registrazione del phantom cervicale}
\cfoot{\thepage}
\setlength{\headheight}{14pt}

\setlist[itemize]{topsep=3pt,itemsep=2pt}
\setlist[enumerate]{topsep=3pt,itemsep=2pt}

\definecolor{todocolor}{RGB}{145,75,0}
\newcommand{\todo}[1]{%
    \par\noindent
    \colorbox{yellow!25}{%
        \parbox{0.94\linewidth}{\textcolor{todocolor}{\textbf{Da completare:} #1}}%
    }\par
}
\newcommand{\R}{\mathbf{R}}
\newcommand{\tvec}{\mathbf{t}}
\newcommand{\x}{\mathbf{x}}
\newcommand{\y}{\mathbf{y}}
\newcommand{\code}[1]{\texttt{\detokenize{#1}}}
\newcommand{\SI}[2]{\ensuremath{#1\,#2}}
\newcommand{\si}[1]{\ensuremath{#1}}
\newcommand{\milli}{\mathrm{m}}
\newcommand{\metre}{\mathrm{m}}
\newcommand{\degree}{^\circ}

\lstset{
    basicstyle=\ttfamily\small,
    frame=single,
    breaklines=true,
    columns=fullflexible,
    showstringspaces=false,
    backgroundcolor=\color{black!2}
}

\title{%
    \textbf{Registrazione learning-based del phantom cervicale\\
    senza punti fiduciali prescritti}\\[0.4em]
    \large Bozza della relazione tecnica e cronologia dello sviluppo
}
\author{Giovanni Contursi \and \textit{Altri componenti del gruppo da inserire}}
\date{Stato del progetto: 22 settembre 2026}

\begin{document}

\maketitle

\begin{abstract}
Questo progetto studia la registrazione rigida tra un insieme parziale e non
ordinato di punti acquisiti con un puntatore ottico e il digital twin di un
phantom cervicale. Il metodo scelto è Deep Closest Point (DCP), che apprende
descrittori e corrispondenze soft e ricava la trasformazione rigida mediante
una testa SVD differenziabile. Sono considerate due modalità di acquisizione:
uno \emph{sweep} denso e connesso e un insieme \emph{sparse} di pochi punti.
Il confronto di riferimento è l'algoritmo clinico-esistente a quattro
landmark con corrispondenze note, implementato mediante Procrustes/SVD.

Il lavoro non si è limitato a eseguire il codice DCP originale. È stato
costruito un generatore sintetico specifico per la mesh del phantom, sono
stati corretti diversi protocolli iniziali troppo favorevoli, è stata resa
possibile la registrazione partial-to-global con cardinalità differenti, sono
state introdotte metriche invarianti all'ordine, una baseline ufficiale, il
calcolo della target registration error (TRE), controlli di riproducibilità e
una supervisione geometrica opzionale sulle corrispondenze. I risultati
attuali mostrano che la loss geometrica migliora nettamente la modalità
sparse, mentre per lo sweep è necessario un peso più moderato. La baseline a
quattro punti rimane molto più accurata, ma richiede landmark noti e un ordine
di acquisizione prescritto; il metodo appreso affronta invece un problema più
difficile e completamente non ordinato.
\end{abstract}

\tableofcontents
\newpage

\section{Contesto e obiettivo}

Il laboratorio dispone di un phantom fisico del distretto cervicale, del suo
modello digitale triangolare e di un sistema di tracking ottico NDI Polaris
Vega XT con puntatore calibrato. Il sistema attuale richiede di toccare quattro
punti noti in un ordine prestabilito. Le quattro corrispondenze permettono di
stimare in forma chiusa la trasformazione rigida dal frame Polaris al frame
del phantom.

L'obiettivo del progetto è rimuovere due vincoli operativi:

\begin{itemize}
    \item non richiedere punti anatomici prescritti;
    \item non richiedere un ordine di acquisizione noto.
\end{itemize}

L'operatore deve poter acquisire un insieme parziale e non ordinato di punti
sulla superficie raggiungibile. La rete deve stimare la stessa trasformazione

\begin{equation}
    \x_{\mathrm{phantom}} = \R\x_{\mathrm{Polaris}} + \tvec,
    \qquad \R\in SO(3),\quad \tvec\in\mathbb{R}^3.
\end{equation}

Tra le varianti proposte dalla traccia è stata scelta V2, Deep Closest Point
(DCP), perché mantiene la soluzione SVD/Procrustes del metodo esistente ma
sostituisce le corrispondenze imposte con corrispondenze apprese. Il confronto
concettuale è quindi particolarmente diretto.

\subsection{Modalità di acquisizione}

Sono state implementate entrambe le modalità richieste:

\begin{description}
    \item[Sweep.] Una patch superficiale connessa, densa e locale. Nel
    protocollo iniziale vengono usati 512 punti e un raggio geodetico-euclideo
    di circa \SI{20}{\milli\metre} attorno alla faccia seme.
    \item[Sparse.] Un insieme non ordinato di punti sparsi sulla superficie.
    Nel protocollo iniziale vengono usati 25 punti, all'interno dell'intervallo
    20--50 indicato dalla traccia.
\end{description}

\todo{Verificare nella versione finale della relazione la terminologia della
superficie acquisibile. La traccia parla di anatomia interna esposta, mentre
un successivo chiarimento della docente ha indicato la superficie esterna
(pelle). Gli esperimenti attuali usano \code{Prototipo_Scocca Esterna.stl},
rinominata \code{phantom_surface.stl}. La scelta va documentata con il
chiarimento ufficiale più recente.}

\section{Fondamenti teorici}

\subsection{Registrazione rigida con corrispondenze note}

Dati due insiemi di punti corrispondenti $\{\x_i\}$ e $\{\y_i\}$, il problema
classico consiste nel minimizzare

\begin{equation}
    \min_{\R,\tvec}
    \sum_{i=1}^{N}\left\|\y_i-(\R\x_i+\tvec)\right\|_2^2,
    \qquad \R^\top\R=\mathbf{I},\quad \det(\R)=1.
\end{equation}

Sottratti i centroidi, si costruisce la matrice di covarianza
$\mathbf{H}=\sum_i(\x_i-\bar{\x})(\y_i-\bar{\y})^\top$. Dalla decomposizione
$\mathbf{H}=\mathbf{U}\mathbf{\Sigma}\mathbf{V}^\top$ si ottiene

\begin{equation}
    \R=\mathbf{V}\mathbf{D}\mathbf{U}^\top,
    \qquad
    \tvec=\bar{\y}-\R\bar{\x},
\end{equation}

dove $\mathbf{D}=\operatorname{diag}(1,1,d)$ corregge le riflessioni e
garantisce $\det(\R)=1$ \cite{arun1987}.

La baseline ufficiale usa esattamente questo algoritmo su quattro landmark
noti. DCP usa lo stesso principio, ma sostituisce ogni target noto con un punto
virtuale ottenuto dalle corrispondenze probabilistiche.

\subsection{Deep Closest Point}

DCP contiene tre blocchi principali \cite{wang2019dcp}:

\begin{enumerate}
    \item un estrattore di feature, nel nostro caso DGCNN;
    \item un Transformer che scambia informazione tra source e target;
    \item una testa SVD che trasforma le corrispondenze soft in $(\R,\tvec)$.
\end{enumerate}

Indichiamo con $\mathbf{F}_s\in\mathbb{R}^{d\times N_s}$ e
$\mathbf{F}_t\in\mathbb{R}^{d\times N_t}$ gli embedding. I punteggi di
corrispondenza sono

\begin{equation}
    \mathbf{S}=\frac{\mathbf{F}_s^\top\mathbf{F}_t}{\sqrt{d}},
    \qquad
    \mathbf{W}=\operatorname{softmax}_{t}(\mathbf{S}).
\end{equation}

Per ogni punto source viene costruito un target virtuale

\begin{equation}
    \widehat{\y}_i=\sum_{j=1}^{N_t}W_{ij}\y_j.
\end{equation}

La SVD è quindi applicata alle coppie $(\x_i,\widehat{\y}_i)$. Questa
formulazione supporta naturalmente $N_s\neq N_t$, condizione necessaria nel
nostro protocollo: il source è parziale, mentre il target è il digital twin
globale.

\subsection{Metriche}

Le metriche principali sono:

\begin{description}
    \item[Errore di rotazione.] Angolo della rotazione relativa
    $\R_{\mathrm{pred}}\R_{\mathrm{gt}}^\top$, espresso in gradi.
    \item[Errore di traslazione.] Norma euclidea
    $\|\tvec_{\mathrm{pred}}-\tvec_{\mathrm{gt}}\|_2$, in millimetri.
    \item[Source proxy.] Disaccordo medio, sui punti osservati, tra la
    trasformazione predetta e la trasformazione ground truth. È una metrica
    diagnostica e non deve essere confusa con la TRE anatomica.
    \item[TRE.] Errore sui target clinicamente significativi. Ogni target,
    definito nel frame phantom, viene trasformato nel frame Polaris con
    l'inversa della ground truth e poi riportato nel phantom con la posa
    stimata:
    \begin{equation}
        e_k=\left\|
        \R_{\mathrm{pred}}
        \R_{\mathrm{gt}}^\top(\mathbf{p}_k-\tvec_{\mathrm{gt}})
        +\tvec_{\mathrm{pred}}-\mathbf{p}_k
        \right\|_2.
    \end{equation}
\end{description}

Per ogni metrica vengono riportati media, deviazione standard, mediana,
95-esimo percentile e massimo. Media e mediana descrivono il caso tipico;
$p_{95}$ e massimo evidenziano le code e i fallimenti.

\section{Evoluzione del progetto}

Questa sezione documenta anche i tentativi che non sono entrati nel protocollo
finale. Sono importanti perché spiegano le scelte metodologiche e impediscono
di interpretare come reali risultati ottenuti su problemi troppo semplici.

\subsection{Analisi iniziale di traccia, articoli e modelli}

Il lavoro è iniziato dall'analisi della traccia, dell'articolo DCP, del metodo
least-squares/SVD e dei componenti STL del phantom. In questa fase sono stati
chiariti:

\begin{itemize}
    \item la direzione della trasformazione, Polaris $\rightarrow$ phantom;
    \item l'uso dei millimetri nel frame fisico;
    \item la necessità di confrontare DCP e baseline sulla stessa TRE;
    \item la separazione obbligatoria tra sweep e sparse;
    \item la necessità di non usare quattro vertici casuali come baseline.
\end{itemize}

Inizialmente mancavano le coordinate ufficiali dei quattro landmark e dei
target TRE. Per evitare una baseline artificiale, tali campi sono rimasti
vuoti fino alla consegna del codice di registrazione e delle coordinate da
parte dei docenti.

\subsection{Adattamento del codice DCP al phantom}

Il codice di riferimento DCP era progettato per dataset generici di point
cloud. È stato introdotto \code{PhantomDataset}, che:

\begin{itemize}
    \item carica e valida la mesh STL con \code{trimesh};
    \item elimina triangoli degeneri e verifica coordinate e aree finite;
    \item campiona uniformemente punti sulle facce usando coordinate
    baricentriche corrette;
    \item costruisce il grafo di adiacenza delle facce per generare patch
    connesse nello sweep;
    \item genera trasformazioni rigide casuali con rotazioni Euler e
    traslazioni limitate;
    \item simula il source nel frame Polaris e aggiunge rumore gaussiano;
    \item normalizza le coordinate fisiche mediante una scala comune di
    \SI{100}{\milli\metre}.
\end{itemize}

Il test geometrico standalone verifica che la ground truth sia internamente
coerente e visualizza la posizione della patch sulla mesh.

\subsection{Prima versione: corrispondenze implicite troppo facili}

Nella prima implementazione source e target contenevano gli stessi punti, nello
stesso ordine, prima e dopo la trasformazione. Il modello poteva quindi
sfruttare una corrispondenza uno-a-uno nascosta. Un test di overfitting produsse
circa \SI{0.20}{\degree} di errore di rotazione e
\SI{0.39}{\milli\metre} di errore di traslazione, ma il dato non rappresentava
il problema operativo richiesto.

La conseguenza è metodologicamente importante: un errore basso non è
sufficiente a dimostrare la correttezza del protocollo. È necessario che il
dataset non fornisca accidentalmente alla rete l'informazione che dovrebbe
imparare.

\subsection{Permutazioni e metriche invarianti all'ordine}

Source e target sono stati successivamente permutati in modo indipendente.
Questo ha fatto emergere un secondo problema: le metriche MSE/MAE iniziali
confrontavano direttamente
\code{transformed_src} con \code{target}. Tale confronto è valido soltanto se
i punti sono corrispondenti e ordinati.

Le metriche di training e validation sono state corrette confrontando la
trasformazione predetta e la ground truth applicate allo stesso insieme di
punti. Per il verso diretto, ad esempio,

\begin{equation}
    \operatorname{MSE}_{AB}=
    \frac{1}{3N_s}\sum_i
    \left\|
      (\R_{\mathrm{pred}}\x_i+\tvec_{\mathrm{pred}})
      -(\R_{\mathrm{gt}}\x_i+\tvec_{\mathrm{gt}})
    \right\|_2^2.
\end{equation}

In questo modo la metrica misura davvero l'errore di posa e non dipende
dall'ordine o dalla discretizzazione delle point cloud.

\subsection{Campionamenti indipendenti}

È stato poi rimosso il riuso delle stesse coordinate geometriche tra source e
target. Un test di overfitting su point set indipendenti ma relativi alla stessa
regione diede circa \SI{0.35}{\degree} e \SI{0.69}{\milli\metre}. Il risultato
confermava che l'architettura e la SVD erano addestrabili, ma il problema era
ancora locale-locale e non coincideva con la registrazione richiesta.

\subsection{Passaggio definitivo a partial-to-global}

La modifica più importante è stata definire correttamente i due ingressi:

\begin{itemize}
    \item \textbf{source}: acquisizione parziale, 512 punti sweep oppure 25
    punti sparse, espressa nel frame Polaris e affetta da rumore;
    \item \textbf{target}: discretizzazione globale del digital twin, 1024
    punti nel frame phantom.
\end{itemize}

Il target globale viene creato una volta con Farthest Point Sampling (FPS) a
partire da un insieme più ampio di candidati uniformi. Il seed dedicato rende
la discretizzazione deterministica. Il target viene inoltre salvato nel
checkpoint e ripristinato durante la valutazione: training, validation e test
usano così esattamente lo stesso digital twin.

Il passaggio a partial-to-global ha aumentato notevolmente la difficoltà. I
primi smoke test mostrarono errori anche superiori a
\SI{100}{\milli\metre}; ciò non era un semplice peggioramento del modello, ma
la conseguenza del passaggio dal problema artificiale con corrispondenze
implicite al problema reale con overlap parziale e molte ambiguità.

\subsection{Cardinalità diverse nella testa SVD}

La testa SVD originale assumeva implicitamente cardinalità compatibili. È
stata riscritta per produrre una matrice di punteggi
$N_s\times N_t$ e un target virtuale per ciascun punto source. È stata inoltre
aggiunta la correzione esplicita delle riflessioni per garantire
$\det(\R)=+1$. Questa modifica consente, ad esempio, di registrare 25 punti
sparse contro 1024 punti del digital twin.

\subsection{Stabilità con batch piccoli}

I primi esperimenti di overfitting mostrarono una forte differenza tra
training e validation. Una delle cause era l'uso di Batch Normalization con
batch molto piccoli. Nel backbone DGCNN la BatchNorm è stata sostituita con
GroupNorm, scegliendo automaticamente il maggior numero di gruppi compatibile
fino a otto. La GroupNorm non dipende dalle statistiche correnti del batch e
ha reso più coerenti training ed evaluation.

\subsection{Baseline ufficiale e target TRE}

Dopo la consegna del codice di registrazione dei docenti sono stati inseriti i
quattro landmark ufficiali nel frame phantom:

\begin{equation}
\begin{split}
    &(2.5,-65.0,82.5),\quad (50.0,-65.0,77.5),\\
    &(50.0,25.0,50.0),\quad (42.5,102.5,27.5)\quad [\si{\milli\metre}].
\end{split}
\end{equation}

Per ogni campione i landmark vengono portati nel frame Polaris con l'inversa
della ground truth, viene aggiunto lo stesso rumore del tracker e la posa è
stimata con SVD classica. La baseline e DCP sono quindi valutati sugli stessi
campioni e sugli stessi target TRE.

Sono stati definiti cinque target lungo la regione centrale cervicale, con
particolare attenzione alla zona del midollo. La correttezza matematica della
pipeline è verificata con casi sintetici: recupero esatto senza rumore, TRE
nullo, traslazione di \SI{1}{\milli\metre} e rotazione nota di
\SI{90}{\degree}.

\subsection{Supervisione geometrica delle corrispondenze}

La sola pose loss supervisiona l'uscita finale ma non impone direttamente che
la matrice di attenzione rappresenti corrispondenze geometriche sensate. È
stata quindi aggiunta una loss opzionale.

La ground truth allinea ogni punto source nel frame phantom. Le distanze dai
punti del target costruiscono una distribuzione soft

\begin{equation}
    q_{ij}=\operatorname{softmax}_j\left(
      -\frac{\|\R_{gt}\x_i+\tvec_{gt}-\y_j\|^2}{2\sigma^2}
    \right),
\end{equation}

con $\sigma=\SI{5}{\milli\metre}$. La rete produce
$p_{ij}=\operatorname{softmax}_j(S_{ij})$ e viene minimizzata

\begin{equation}
    \mathcal{L}_{corr}=D_{KL}(q\|p).
\end{equation}

L'obiettivo totale è

\begin{equation}
    \mathcal{L}_{tot}=\mathcal{L}_{pose}
    +0.1\mathcal{L}_{cycle}
    +\lambda_{corr}\mathcal{L}_{corr},
\end{equation}

dove il termine di ciclo è presente soltanto se richiesto. La ground truth è
usata esclusivamente per costruire la loss durante l'addestramento e non viene
fornita agli ingressi della rete. Il checkpoint migliore continua a essere
selezionato sulla pose validation loss, senza includere la KL, per non favorire
un modello che predice buone corrispondenze ma una posa peggiore.

\subsection{Riproducibilità e controlli}

Sono stati aggiunti:

\begin{itemize}
    \item seed deterministici per campione e split separati;
    \item salvataggio di argomenti, ottimizzatore, scheduler, epoca, digital
    twin e hash SHA-256 della mesh;
    \item distinzione tra \code{--resume} e \code{--pretrained};
    \item blocco della sovrascrittura accidentale di esperimenti esistenti;
    \item backup dei file sorgente nella directory di ogni esperimento;
    \item clipping del gradiente e controlli su valori non finiti;
    \item otto unit test su forward/backward, scale, permutazioni,
    corrispondenze, checkpoint, dataset, SVD e TRE.
\end{itemize}

Tutti gli otto test risultano superati sia su Windows sia sul computer Linux
usato per i training GPU.

\section{Architettura e protocollo attuale}

\subsection{Configurazione del modello}

La configurazione principale usa:

\begin{table}[h]
\centering
\caption{Configurazione DCP principale.}
\begin{tabular}{ll}
\toprule
Componente & Configurazione \\
\midrule
Backbone & DGCNN, $k=20$ \\
Dimensione embedding & 512 \\
Pointer & Transformer, 1 blocco, 4 teste \\
Feed-forward & 1024 unità \\
Dropout & 0 \\
Testa & SVD differenziabile \\
Parametri addestrabili & 5\,568\,896 \\
Normalizzazione & GroupNorm \\
\bottomrule
\end{tabular}
\end{table}

\subsection{Generazione dei dati}

Per ciascun indice del dataset viene creato un generatore indipendente a
partire da \code{SeedSequence([seed, index])}. Ciò rende ogni campione
riproducibile senza riutilizzare la stessa sequenza casuale.

I parametri finora usati nei training principali sono:

\begin{table}[h]
\centering
\caption{Protocollo sintetico di riferimento.}
\begin{tabular}{ll}
\toprule
Parametro & Valore \\
\midrule
Rotazione massima per asse & \SI{0.17}{rad} ($\approx\SI{9.74}{\degree}$) \\
Traslazione massima per asse & \SI{10}{\milli\metre} \\
Rumore per coordinata & $\sigma=\SI{0.3}{\milli\metre}$ \\
Scala rete & \SI{100}{\milli\metre} \\
Target globale & 1024 punti FPS \\
Sweep iniziale & 512 punti, raggio \SI{20}{\milli\metre} \\
Sparse iniziale & 25 punti \\
Training / validation & 5000 / 500 campioni \\
Epoche & 50 \\
Ottimizzatore & Adam, $10^{-3}$, weight decay $10^{-4}$ \\
Scheduler & StepLR, fattore 0.5 dopo 25 epoche \\
Gradient clipping & norma massima 1 \\
\bottomrule
\end{tabular}
\end{table}

\subsection{Loss di posa}

La loss di rotazione non confronta direttamente gli elementi di due matrici
arbitrarie, ma misura quanto la rotazione relativa si discosta dall'identità:

\begin{equation}
    \mathcal{L}_{R}=\operatorname{MSE}
    (\R_{pred}^{\top}\R_{gt},\mathbf{I}),
    \qquad
    \mathcal{L}_{t}=\operatorname{MSE}
    (\tvec_{pred},\tvec_{gt}).
\end{equation}

Le traslazioni sono espresse nelle unità normalizzate della rete durante il
training e convertite una sola volta in millimetri durante l'evaluation.

\section{Risultati sperimentali disponibili}

\subsection{Come leggere i risultati}

I valori di validation loss sono utili per selezionare checkpoint all'interno
dello stesso protocollo, ma non sono direttamente interpretabili come
millimetri. La metrica primaria del progetto è la TRE in evaluation su campioni
held-out. Gli esperimenti iniziali con source e target corrispondenti sono
considerati test di debug e non risultati finali.

\subsection{Ablation del peso della loss geometrica}

Il peso $\lambda_{corr}$ è stato confrontato sullo stesso development set di
2000 campioni, seed 44. In questo modo la differenza è attribuibile al modello
e non a campioni casuali differenti.

\begin{table}[ht]
\centering
\caption{Ablation di $\lambda_{corr}$ in modalità sweep, development seed 44.
Tutti gli errori lineari sono in millimetri.}
\resizebox{\textwidth}{!}{%
\begin{tabular}{crrrrrrr}
\toprule
$\lambda_{corr}$ & {Rot. media [deg]} & {Trasl. media} & {Proxy medio}
& {TRE media} & {TRE mediana} & {TRE $p_{95}$} & {TRE max} \\
\midrule
0.00 & 7.963 & 12.971 & 17.801 & 14.028 & 13.390 & 24.843 & 45.863 \\
0.01 & 8.067 & 13.351 & 17.281 & 14.061 & 13.258 & 25.815 & 52.394 \\
0.03 & 7.631 & 12.504 & 13.931 & \bfseries 12.761 & \bfseries 11.582 & \bfseries 24.660 & 57.856 \\
0.10 & 7.771 & 13.421 & \bfseries 11.991 & 13.433 & 11.903 & 27.260 & 90.501 \\
\bottomrule
\end{tabular}}
\end{table}

Per lo sweep, $\lambda_{corr}=0.03$ è il compromesso migliore. Rispetto a
$\lambda=0$, la TRE media scende da 14.028 a
\SI{12.761}{\milli\metre}, un miglioramento di circa il 9\%, e la mediana
migliora di circa il 13.5\%. Il peso 0.1 ottiene il proxy migliore ma peggiora
media, $p_{95}$ e soprattutto il massimo. Ciò indica che una supervisione
troppo forte può ottimizzare le corrispondenze locali senza garantire la
robustezza della posa globale.

\begin{table}[ht]
\centering
\caption{Ablation di $\lambda_{corr}$ in modalità sparse, development seed 44.
Tutti gli errori lineari sono in millimetri.}
\resizebox{\textwidth}{!}{%
\begin{tabular}{crrrrrrr}
\toprule
$\lambda_{corr}$ & {Rot. media [deg]} & {Trasl. media} & {Proxy medio}
& {TRE media} & {TRE mediana} & {TRE $p_{95}$} & {TRE max} \\
\midrule
0.00 & 3.947 & 5.485 & 8.333 & 5.440 & 5.152 & 9.539 & \bfseries 15.988 \\
0.01 & 3.394 & 4.270 & 6.585 & 4.061 & 3.724 & 7.739 & 18.519 \\
0.03 & 3.174 & 3.996 & 5.994 & 3.653 & 3.321 & 7.105 & 18.310 \\
0.10 & \bfseries 3.073 & \bfseries 3.899 & \bfseries 5.838
& \bfseries 3.605 & \bfseries 3.226 & \bfseries 7.062 & 17.142 \\
\bottomrule
\end{tabular}}
\end{table}

Per la modalità sparse il comportamento è più netto: aumentando il peso fino
a 0.1 migliorano tutte le statistiche aggregate. Rispetto a $\lambda=0$, la
TRE media migliora di circa il 33.7\%, la mediana del 37.4\% e il $p_{95}$ del
26.0\%. Il massimo resta leggermente peggiore del modello senza KL, ma 0.1 è
anche il migliore tra tutti i modelli supervisionati nelle code.

Le configurazioni provvisoriamente selezionate sono quindi:

\begin{itemize}
    \item sweep: $\lambda_{corr}=0.03$;
    \item sparse: $\lambda_{corr}=0.10$.
\end{itemize}

\subsection{Confronto con la baseline a quattro punti}

Sul development set seed 44, la baseline presenta una TRE media di circa
\SI{0.40}{\milli\metre} in entrambe le modalità. Tale risultato è coerente con
il vantaggio di conoscere quattro corrispondenze esatte e ben distribuite. Il
confronto sintetico corrente è riassunto in Tabella~\ref{tab:best-dev}.

\begin{table}[ht]
\centering
\caption{Migliori configurazioni correnti sul development set seed 44.}
\label{tab:best-dev}
\begin{tabular}{llrrrr}
\toprule
Modalità & Metodo & Rot. [deg] & Trasl. [mm] & TRE [mm] & TRE $p_{95}$ [mm] \\
\midrule
Sweep & DCP, $\lambda=0.03$ & 7.631 & 12.504 & 12.761 & 24.660 \\
Sweep & Baseline 4 punti & 0.435 & 0.561 & 0.401 & 0.792 \\
Sparse & DCP, $\lambda=0.10$ & 3.073 & 3.899 & 3.605 & 7.062 \\
Sparse & Baseline 4 punti & 0.434 & 0.564 & 0.401 & 0.795 \\
\bottomrule
\end{tabular}
\end{table}

La baseline è dunque ancora nettamente più accurata. Tuttavia i due metodi non
richiedono lo stesso intervento: la baseline usa soltanto quattro punti ma con
identità e ordine noti; DCP usa punti arbitrari e non ordinati. Il risultato va
presentato senza sostenere che DCP abbia già sostituito la baseline in termini
di accuratezza. Il contributo attuale è la rimozione dei vincoli di
acquisizione, con un costo ancora significativo in TRE.

\subsection{Perché sparse è attualmente migliore di sweep}

A prima vista può sembrare sorprendente che 25 punti sparse superino uno sweep
di 512 punti. Nel generatore corrente, però, i punti sparse sono distribuiti
sull'intera superficie, mentre lo sweep occupa una singola patch locale. Una
patch della scocca può essere liscia, quasi simmetrica o geometricamente poco
distintiva. Molti punti nella stessa zona non aggiungono necessariamente
informazione globale. Pochi punti distanti tra loro possono invece vincolare
meglio rotazione e traslazione.

Questa osservazione ha motivato l'ablation attualmente in esecuzione su:

\begin{itemize}
    \item raggio dello sweep: 20, 40 e \SI{60}{\milli\metre};
    \item numero di punti sparse: 20, 25, 35 e 50;
    \item confronto, per ogni acquisizione, tra pose loss standard e migliore
    loss geometrica selezionata per la modalità.
\end{itemize}

\todo{Inserire qui la tabella dell'ablation di acquisizione al termine del
training \code{acq_sweep_r40_dcp_seed42} e degli esperimenti successivi. Non
selezionare il protocollo finale prima di aver confrontato tutti i risultati
sullo stesso development seed.}

\section{Problemi incontrati e conseguenze}

\begin{longtable}{p{0.23\textwidth}p{0.34\textwidth}p{0.34\textwidth}}
\caption{Principali problemi tecnici e metodologici affrontati.}\\
\toprule
Problema & Correzione & Conseguenza \\
\midrule
\endfirsthead
\toprule
Problema & Correzione & Conseguenza \\
\midrule
\endhead
Percorso STL errato & Individuazione dei file reali in
\code{ModelloPhantom} e uso di percorsi quotati & Pipeline eseguibile su
Windows e Linux; scelta esplicita della mesh. \\

Baseline con vertici casuali & Attesa delle coordinate ufficiali e
implementazione del metodo SVD dei docenti & Confronto difendibile e coerente
con il sistema esistente. \\

Source e target identici & Permutazioni, campionamenti indipendenti e infine
target globale FPS & Eliminazione delle corrispondenze implicite; problema più
realistico ma più difficile. \\

Metriche pointwise non valide & Confronto tra trasformazione predetta e GT
applicate agli stessi punti & MSE e MAE invarianti all'ordine. \\

Cardinalità source/target diverse & Matrice di attenzione $N_s\times N_t$ e
target virtuali source-to-target & Supporto per 25/512 punti contro 1024. \\

Gap train/eval con batch piccoli & BatchNorm sostituita da GroupNorm &
Comportamento più stabile e checkpoint più affidabili. \\

Digital twin ricampionato & Target FPS fisso salvato nel checkpoint &
Riproducibilità tra training, validation e test. \\

Unità mm / unità rete & Normalizzazione unica e conversione esplicita in
evaluation & Errori fisici interpretabili e riduzione dei rischi di fattori
100 o 1000. \\

Attenzione non supervisionata & KL geometrica con etichette soft & Grande
miglioramento sparse; peso da calibrare nello sweep. \\

Esperimenti sovrascrivibili & Blocco su directory esistenti, backup sorgenti,
checkpoint best/last & Maggiore tracciabilità; le cartelle incomplete devono
essere riprese o archiviate esplicitamente. \\
\bottomrule
\end{longtable}

\section{Validazione sperimentale ancora necessaria}

Prima di considerare chiusa la parte sperimentale sono previste le seguenti
attività.

\subsection{Ablation dell'acquisizione}

L'esperimento in corso deve determinare se le prestazioni cambiano con una
patch sweep più estesa e con un numero maggiore di punti sparse. I parametri
devono essere scelti sul development set, non sul test finale.

\subsection{Ripetizione su più seed di training}

Una sola inizializzazione non consente di distinguere una scelta robusta da un
risultato fortunato. Dopo aver fissato raggio e numero di punti, le quattro
configurazioni principali dovranno essere ripetute con almeno tre seed di
training:

\begin{itemize}
    \item sweep senza loss geometrica;
    \item sweep con $\lambda_{corr}=0.03$;
    \item sparse senza loss geometrica;
    \item sparse con $\lambda_{corr}=0.10$.
\end{itemize}

La relazione finale dovrà riportare media e deviazione tra seed, oltre alle
statistiche per campione.

\subsection{Test finale held-out}

Il seed finale deve rimanere inutilizzato durante le scelte progettuali. La
procedura proposta usa:

\begin{itemize}
    \item seed 44 per sviluppo e ablation;
    \item seed di training distinti, ad esempio 42, 123 e 777;
    \item seed 314159 per la sola valutazione finale.
\end{itemize}

\subsection{Bacino di convergenza e success rate}

La traccia richiede di studiare il disallineamento iniziale entro il quale il
metodo converge. Occorre valutare più livelli di rotazione e traslazione, ad
esempio aumentando separatamente i limiti, e definire prima dell'esperimento
una soglia di successo basata sulla TRE. Un esempio da discutere è
$\mathrm{TRE}<\SI{5}{\milli\metre}$, ma la soglia finale dovrebbe essere
motivata clinicamente o concordata con la docente.

\subsection{Robustezza al rumore}

È opportuno confrontare almeno $\sigma=0$, 0.3 e
\SI{0.6}{\milli\metre}. Questo permette di separare l'errore dovuto
all'ambiguità geometrica da quello dovuto al tracker.

\subsection{Reality check con Polaris}

Il limite principale è l'assenza, al momento, di una piccola acquisizione
reale sul phantom. La traccia richiede un controllo finale con Polaris. I
risultati sintetici non dimostrano da soli la generalizzazione a errori reali
di calibrazione, contatto, occlusione e campionamento manuale.

\section{Struttura del codice}

\begin{description}
    \item[\code{phantom_data.py}] Caricamento mesh, validazione, campionamento
    uniforme, grafo delle facce, patch sweep, FPS, trasformazioni casuali e
    dataset partial-to-global.
    \item[\code{model.py}] DGCNN/PointNet, Transformer, testa SVD e modello
    DCP. Contiene GroupNorm e supporto opzionale alla restituzione dei logits
    di corrispondenza.
    \item[\code{correspondence_loss.py}] Costruzione delle etichette
    geometriche soft e KL $D_{KL}(q\|p)$.
    \item[\code{train.py}] Training e validation, loss di posa/ciclo/KL,
    checkpoint, backup dei sorgenti, resume/pretrained, logging e TensorBoard.
    \item[\code{evaluate.py}] Ricostruzione della configurazione dal
    checkpoint, metriche fisiche, TRE, baseline ufficiale e riepilogo
    statistico.
    \item[\code{test_correspondence.py}] Otto test automatici della pipeline.
    \item[\code{util.py}] Trasformazioni di point cloud e conversioni di
    rotazione.
\end{description}

\section{Conclusioni provvisorie}

Il progetto ha raggiunto una pipeline completa e verificata per la
registrazione sintetica partial-to-global del phantom cervicale. Il progresso
più rilevante non è un singolo valore di loss, ma il passaggio da un test
iniziale con corrispondenze implicite a un protocollo realistico, non ordinato,
con overlap parziale, digital twin globale fisso, baseline ufficiale e TRE
anatomica.

I risultati attuali supportano tre conclusioni:

\begin{enumerate}
    \item la supervisione geometrica delle corrispondenze è utile, soprattutto
    nella modalità sparse;
    \item uno sweep locale molto denso può essere meno informativo di pochi
    punti distribuiti globalmente;
    \item la baseline a quattro landmark rimane molto più accurata, ma paga
    tale accuratezza con corrispondenze e ordine prescritti.
\end{enumerate}

Il progetto non deve quindi essere presentato come una vittoria numerica di
DCP sulla baseline. La domanda scientifica corretta è quanto errore si accetta
per eliminare la procedura vincolata dei quattro landmark e quali modalità di
acquisizione rendono questo compromesso più favorevole. Le prossime ablation,
la valutazione multi-seed e il reality check Polaris serviranno a rispondere in
modo più solido.

\appendix

\section{Cronologia sintetica}

\begin{longtable}{p{0.17\textwidth}p{0.75\textwidth}}
\toprule
Fase & Attività principale \\
\midrule
\endhead
Analisi & Studio della traccia, DCP, soluzione SVD e componenti STL. \\
Dataset V1 & Generazione di coppie sintetiche da mesh; verifica di frame,
scala e rumore. \\
Debug & Overfitting su pochi campioni e controllo che la SVD fosse
differenziabile e stabile. \\
Unordered & Permutazione dei punti e correzione delle metriche non invarianti
all'ordine. \\
Independent & Rimozione del riuso delle stesse coordinate tra point set. \\
Partial-to-global & Source parziale contro digital twin globale FPS da 1024
punti. \\
Stabilità & GroupNorm, gradient clipping e controlli sui valori non finiti. \\
Baseline & Inserimento dei quattro landmark ufficiali e SVD classica. \\
TRE & Inserimento dei target, metriche per target e test matematici. \\
Supervisione & Loss KL geometrica sulle corrispondenze source-to-target. \\
Ablation loss & Confronto $\lambda\in\{0,0.01,0.03,0.10\}$ sul seed 44. \\
Ablation acquisizione & Confronto in corso di raggi sweep e cardinalità sparse. \\
Fase finale & Multi-seed, robustezza, basin of attraction, success rate e dati
Polaris reali. \\
\bottomrule
\end{longtable}

\section{Esempio di comando riproducibile}

Il seguente comando mostra la configurazione base di un training sweep con
supervisione geometrica. I valori definitivi di raggio e peso dovranno essere
aggiornati dopo le ablation.

\begin{lstlisting}[language=bash]
python train.py \
  --stl phantom_surface.stl \
  --exp_name sweep_corr003_seed42 \
  --emb_nn dgcnn --pointer transformer --head svd \
  --emb_dims 512 --ff_dims 1024 \
  --n_blocks 1 --n_heads 4 --dgcnn_k 20 \
  --batch_size 2 --epochs 50 --scheduler 25 \
  --lr 0.001 --grad_clip 1 \
  --dset_mode sweep \
  --dset_num_samples 5000 --val_num_samples 500 \
  --dset_n_points 512 --target_n_points 1024 \
  --dset_rot_max 0.17 --dset_trans_max 10 \
  --noise_sigma 0.3 --network_scale_mm 100 \
  --patch_radius_mm 20 --manualSeed 42 \
  --num_workers 4 --corr_weight 0.03 \
  --corr_sigma_mm 5
\end{lstlisting}

\begin{thebibliography}{9}

\bibitem{arun1987}
K. S. Arun, T. S. Huang, S. D. Blostein,
``Least-Squares Fitting of Two 3-D Point Sets,''
\emph{IEEE Transactions on Pattern Analysis and Machine Intelligence},
vol. PAMI-9, no. 5, 1987.

\bibitem{wang2019dcp}
Y. Wang, J. M. Solomon,
``Deep Closest Point: Learning Representations for Point Cloud Registration,''
\emph{Proceedings of ICCV}, 2019.

\bibitem{besl1992}
P. J. Besl, N. D. McKay,
``A Method for Registration of 3-D Shapes,''
\emph{IEEE Transactions on Pattern Analysis and Machine Intelligence},
vol. 14, no. 2, 1992.

\end{thebibliography}

\end{document}
